% Paper 3 — Replicable Pipeline Template
% Section 7: How to build a local-language economic narrative index

\section{Replicable Pipeline Template}\label{sec:template}

This section provides a step-by-step guide for researchers who want to build a local-language economic narrative index in another language or country context. The template is designed around a single binding constraint: \textbf{the annotation budget}, measured in both time and cost.

\subsection{Step-by-Step Protocol}

\begin{enumerate}
    \item \textbf{Select a corpus.} Identify a publicly available or scrapable news corpus in the target language with at least 50,000 articles and publication dates. Sources: academic repositories (Mendeley, Zenodo), news API archives, or custom scrapers (Newspaper3k, Scrapy). \textit{Time: 1--2 weeks. Cost: \$0 if using existing corpus; \$50--\$200 for API access.}

    \item \textbf{Define the classification task.} Economic relevance is a binary label (Economic / Not Economic). For more granular tasks (sentiment, topic, narrative frame), adapt the six-field schema from Appendix~\ref{app:annotation-prompt}. Pre-register the annotation protocol and evaluation metrics (primary: F1 for the target class; secondary: Cohen's $\kappa$). \textit{Time: 1 day.}

    \item \textbf{Generate weak labels.} Create keyword-derived pseudo-labels using a domain-specific term list (see Appendix~\ref{app:keywords} for the Bangla economics term list). Expect precision of 5--45\% depending on keyword quality. This step requires no annotation budget and provides training data for the initial classifier. \textit{Time: 2--3 days. Cost: \$0.}

    \item \textbf{Train a baseline classifier.} Use TF-IDF + Logistic Regression (no GPU required). This establishes the reproducibility baseline and reveals the weak-label ceiling before any annotation budget is committed. \textit{Time: 30 minutes. Cost: \$0. Compute: 1 CPU core.}

    \item \textbf{Generate silver-standard labels.} Use an LLM annotator (Claude, GPT-4o) at temperature 0.0 with a structured prompt to annotate 300--500 articles. Self-consistency check on 50 articles. Total API cost at current rates: \$1--\$5 for 300 articles. \textit{Time: 1--2 days. Cost: \$1--\$5. Compute: None (API).}

    \item \textbf{Evaluate the weak-label ceiling.} Compare the TF-IDF baseline against the silver standard. Compute McNemar's test between any two weak-label models to determine whether architectural improvements extract additional signal. If $p > 0.05$, the weak labels are the binding constraint. \textit{Time: 1 hour.}

    \item \textbf{Decide on annotation investment} (see Section~\ref{sec:investment-decision} below). If the weak-label F1 is below 0.5 and resources permit, invest in designed sampling (uncertainty-guided, not random) for the next annotation batch. \textit{Time: varies.}

    \item \textbf{Build the index.} Aggregate monthly scores, compare article-weighted and source-balanced specifications, preserve model and release versions, and validate against available macroeconomic indicators. \textit{Time: 1--2 days after predictions are complete.}
\end{enumerate}

\subsection{Cost Table}\label{sec:cost-table}

Table~\ref{tab:costs} summarizes the compute and financial costs associated with each step, using the BENI pipeline as a reference.

\begin{table}[h]
\centering
\small
\begin{tabular}{lccc}
\toprule
Component & Compute & Time & Cost (USD) \\
\midrule
Corpus acquisition & 1 CPU, 4 GB RAM & 1--2 weeks & \$0--\$200 \\
Keyword weak labels & 1 CPU & 2--3 days & \$0 \\
TF-IDF + LogReg training & 1 CPU, 2 GB RAM & 2 minutes & \$0 (free tier) \\
LLM annotation (300 articles) & API only & 1--2 days & \$1--\$5 \\
Self-consistency check (50 articles) & API only & 1 hour & \$0.20--\$1 \\
Transformer fine-tuning (100k--1M articles) & 16 GB GPU (T4) & 2--8 hours & \$0 (Kaggle) or \$2--\$10 \\
McNemar + evaluation & 1 CPU & 1 minute & \$0 \\
Active learning simulation & 1 CPU, 4 GB RAM & 30 minutes & \$0 \\
Index construction and validation & 1 CPU & 10--60 minutes & \$0 \\
\bottomrule
\end{tabular}
\caption{Compute and cost estimates for each pipeline step. Total: \$3--\$208, primarily corpus-dependent.}
\label{tab:costs}
\end{table}

The total pipeline cost---excluding corpus acquisition---is **under \$10** for the LLM annotation (300 articles) and can be as low as **\$3** when free API credits or open-source models are used. The compute requirements are met by any consumer laptop with 8 GB RAM for the TF-IDF stages; only the BanglaBERT fine-tuning step requires a GPU (available at no cost via Kaggle or Google Colab).

\subsection{Annotation Investment Decision Tree}\label{sec:investment-decision}

Figure~\ref{fig:decision-tree} presents a decision tree for determining the minimum viable annotation investment based on the target base rate and available resources. The key thresholds, derived from the BENI experiments, are:

\begin{itemize}
    \item \textbf{Base rate $\geq$ 20\%}: Random sampling of 200--300 articles with TF-IDF + LogReg achieves F1 $\geq$ 0.6. LLM annotation is optional.
    \item \textbf{Base rate 10--20\%}: Active learning with uncertainty sampling reaches 90\% of full performance at $k \approx 250$. LLM annotation recommended.
    \item \textbf{Base rate $<$ 10\%} (BENI case): Random sampling fails at all budgets up to $k = 300$. Designed sampling (uncertainty-guided, positive-unlabeled learning) or LLM annotation is essential. Minimum budget: $k \geq 500$ for TF-IDF, or $k \geq 200$ with LLM-generated labels.
\end{itemize}

\begin{figure}[h]
\centering
\begin{tabular}{ll}
\toprule
Situation & Recommended Action \\
\midrule
No budget, no API access & Keyword weak labels + TF-IDF baseline only \\
\$5--\$10 budget, base rate $<$ 10\% & LLM annotate 300 articles, TF-IDF evaluation \\
\$10--\$50 budget & LLM annotate 300--500 + designed sampling \\
GPU available (free tier) & Add BanglaBERT fine-tuning on weak labels \\
Human annotators available & 2 annotators $\times$ 300 articles + adjudication \\
\bottomrule
\end{tabular}
\caption{Decision matrix for annotation strategy by budget and setting.}
\label{fig:decision-tree}
\end{figure}

\subsection{Validation Checklist}

Before publishing an index, verify each item on the following checklist:

\begin{enumerate}
    \item \textbf{Annotation reliability}: Report Cohen's $\kappa$, observed agreement, and prevalence-adjusted $\kappa$ for the annotation set.
    \item \textbf{Model ceiling}: Compare against a TF-IDF baseline on identical training labels. Report McNemar's test.
    \item \textbf{Active learning curve}: Report F1 (target class) vs. annotation budget across at least 5 budget levels with confidence intervals.
    \item \textbf{Temporal validation}: Use temporal train/val/test splits (not random). Report performance separately for each period.
    \item \textbf{Macro correlation}: Report both level and first-differenced correlations with at least two independent macroeconomic indicators.
    \item \textbf{Source balance}: Report article-weighted and source-balanced indices, especially when source volume is uneven.
    \item \textbf{Reproducibility}: Publish all code, annotation prompts, and preprocessed data files.
\end{enumerate}
